
\thispagestyle{empty}

\begin{center}
    \small
    \sc Instituto Federal de Educação, Ciência e Tecnologia de Minas Gerais \\ 
    \sc Campus Formiga \\
    \sc Mestrado Profissional em Administração
    
   
    \vspace{3cm}
    \large {Rodrigo Emiliano dos Santos}
    
    \vspace{3cm}
     \large \textbf{Value Investing no Mercado Brasileiro} \\
   Aplicação da Fórmula de Valor Intrínseco de Benjamin Graham na Seleção de Carteiras no Mercado Brasileiro
    
    \vspace{2cm}
    \begin{flushright}
    \begin{minipage}{0.6\textwidth}
    \small
    Projeto apresentado ao Programa de Pós-
    Graduação em Administração do Instituto
    Federal de Educação, Ciência e
    Tecnologia de Minas Gerais - \textit{Campus} Formiga 
    (IFMG - \textit{Campus} Formiga).
    \end{minipage}
    \end{flushright}
    
    \vspace{0.5cm}
    \begin{flushright}
    \small
    Orientador: Prof. Washington Santos Silva.\\
    Coorientador: Prof. Dr. Lelis Pedro de Andrade.
    
    \vspace{0.5cm}
    Linha de Pesquisa: Finanças Corporativas e Investimentos.
    \end{flushright}
    
    \vfill
    Formiga, Minas Gerais \\
    2026
\end{center}


\newpage

\thispagestyle{empty}

\newenvironment{meuresumo}{
  \clearpage
  \small
  \vspace{-1cm}
  \begin{center}
    \bfseries RESUMO
    \vspace{0.5em}
  \end{center}
  \begin{quote}
}{
  \end{quote}
  \vspace{-1.1em}
  \begin{center}
  \begin{minipage}{0.87\textwidth} 
  \textbf{Palavras-chave:} value investing, valor intrínseco, Benjamin Graham.
  \end{minipage}
  \end{center}
  \clearpage
}

\begin{meuresumo}
O desempenho superior das estratégias de investimento em valor tem sido amplamente documentado na literatura financeira internacional. Embora Fama e French (1992) não adotem explicitamente o arcabouço original de Benjamin Graham, seus modelos multifatoriais de apreçamento de ativos fornecem uma explicação robusta para os retornos excedentes persistentes observados em carteiras compostas por ações subavaliadas. Em particular, a compensação por exposições a riscos sistemáticos — como tamanho e valor — oferece uma fundamentação econômica coerente com os princípios subjacentes ao value investing.

Esta dissertação examina o desempenho de carteiras de ações construídas no mercado acionário brasileiro no período de 2015 a 2025. A seleção dos ativos segue os critérios fundamentalistas de Graham, empiricamente adaptados à realidade brasileira em consonância com Palazzo (2018), acrescidos do requisito de uma Margem de Segurança mínima de 33\% em relação ao valor intrínseco estimado. O desempenho das carteiras é avaliado por meio do Capital Asset Pricing Model (CAPM) e de uma sequência de modelos multifatoriais de apreçamento de ativos, incluindo o modelo de três fatores de Fama–French, o modelo de quatro fatores de Carhart e um modelo de cinco fatores baseado nos fatores de risco do NEFIN, contemplando mercado, tamanho, valor, momentum e investimento.

Os resultados empíricos indicam que as carteiras baseadas nos critérios de Graham apresentam uma exposição defensiva às flutuações do mercado e uma sensibilidade pronunciada ao fator tamanho (SMB), sugerindo que a estratégia está fortemente associada ao risco de empresas de menor capitalização no mercado brasileiro. Embora as carteiras apresentem desempenho superior em termos de retornos acumulados e métricas ajustadas ao risco, como o Índice de Sharpe modificado, a ausência de alfas estatisticamente significativos nas regressões multifatoriais indica que tais retornos são, em grande medida, atribuíveis a exposições a riscos sistemáticos, e não à existência de ineficiências persistentes de mercado.

De forma geral, este estudo contribui para a literatura ao integrar os princípios clássicos do value investing, o critério de Margem de Segurança e modelos contemporâneos de apreçamento multifatorial no contexto brasileiro. Os resultados reafirmam a relevância do arcabouço de Graham, ao mesmo tempo em que esclarecem as fontes subjacentes de retorno das carteiras de valor, demonstrando que seu desempenho superior é predominantemente impulsionado pela compensação por riscos sistemáticos, e não por retornos anormais não explicados por fatores de risco consolidados.

\end{meuresumo}




% formatação do título de cada seção
\makeatletter
\def\@makechapterhead#1{%
  \vspace*{50\p@}%
  {\parindent \z@ \raggedright \normalfont
    \ifnum \c@secnumdepth >\m@ne
      \huge\bfseries \thechapter\space
    \fi
    \huge \bfseries #1\par\nobreak
    \vskip 40\p@
  }}
\makeatother


